\documentclass[12pt]{article}

\usepackage[utf8]{inputenc}
\usepackage[T1]{fontenc}
\usepackage{geometry}
\usepackage{amsmath, amssymb}
\usepackage{enumitem}
\usepackage{setspace}

\geometry{margin=1in}
\setlength{\parindent}{0pt}
\setstretch{1.1}

\begin{document}

% --------------------------------------------------------------------
% Cover Page
% --------------------------------------------------------------------

\begin{center}
\begin{tabular}{p{0.45\linewidth}p{0.45\linewidth}}
Name:\hrulefill & Student ID:\hrulefill \\
\end{tabular}

\vspace{2em}

\Large \textbf{ECON 308 -- International Economic Relations} \\[0.5em]
\large \textbf{Final Exam} \\[1em]

\normalsize EXAMINER: \textbf{Muhebullah Karimzada} \\[0.5em]
December 11, 2025
\end{center}



\hrule
\vspace{2em}
\textbf{Instructions:}
\begin{itemize}
\item You are allowed \textbf{120 minutes} to complete the exam.
    \item Show all work. Unsupported answers receive little credit.

    \item Unless stated otherwise, all variables are real and differentiable as needed.
\end{itemize}

\hrule

\vspace{1em}

\textbf{PLEASE: PUT YOUR NAME AND STUDENT ID ON EVERY SHEET!!}

\vspace{2em}

\textbf{Multiple Choice Answers:}

\vspace{0.8em}


\setlength{\tabcolsep}{18pt}

\begin{tabular}{l l l l}
1.\ \rule{1.2cm}{0.15mm} \hspace{1cm} & 
15.\ \rule{1.2cm}{0.15mm} \hspace{1cm} & 
29.\ \rule{1.2cm}{0.15mm} \hspace{1cm} & 
43.\ \rule{1.2cm}{0.15mm} \\[6pt]

2.\ \rule{1.2cm}{0.15mm} & 
16.\ \rule{1.2cm}{0.15mm} & 
30.\ \rule{1.2cm}{0.15mm} & 
44.\ \rule{1.2cm}{0.15mm} \\[6pt]

3.\ \rule{1.2cm}{0.15mm} & 
17.\ \rule{1.2cm}{0.15mm} & 
31.\ \rule{1.2cm}{0.15mm} & 
45.\ \rule{1.2cm}{0.15mm} \\[6pt]

4.\ \rule{1.2cm}{0.15mm} & 
18.\ \rule{1.2cm}{0.15mm} & 
32.\ \rule{1.2cm}{0.15mm} & \\[6pt]

5.\ \rule{1.2cm}{0.15mm} & 
19.\ \rule{1.2cm}{0.15mm} & 
33.\ \rule{1.2cm}{0.15mm} & \\[6pt]

6.\ \rule{1.2cm}{0.15mm} & 
20.\ \rule{1.2cm}{0.15mm} & 
34.\ \rule{1.2cm}{0.15mm} & \\[6pt]

7.\ \rule{1.2cm}{0.15mm} & 
21.\ \rule{1.2cm}{0.15mm} & 
35.\ \rule{1.2cm}{0.15mm} & \\[6pt]

8.\ \rule{1.2cm}{0.15mm} & 
22.\ \rule{1.2cm}{0.15mm} & 
36.\ \rule{1.2cm}{0.15mm} & \\[6pt]

9.\ \rule{1.2cm}{0.15mm} & 
23.\ \rule{1.2cm}{0.15mm} & 
37.\ \rule{1.2cm}{0.15mm} & \\[6pt]

10.\ \rule{1.2cm}{0.15mm} & 
24.\ \rule{1.2cm}{0.15mm} & 
38.\ \rule{1.2cm}{0.15mm} & \\[6pt]

11.\ \rule{1.2cm}{0.15mm} & 
25.\ \rule{1.2cm}{0.15mm} & 
39.\ \rule{1.2cm}{0.15mm} & \\[6pt]

12.\ \rule{1.2cm}{0.15mm} & 
26.\ \rule{1.2cm}{0.15mm} & 
40.\ \rule{1.2cm}{0.15mm} & \\[6pt]

13.\ \rule{1.2cm}{0.15mm} & 
27.\ \rule{1.2cm}{0.15mm} & 
41.\ \rule{1.2cm}{0.15mm} & \\[6pt]

14.\ \rule{1.2cm}{0.15mm} & 
28.\ \rule{1.2cm}{0.15mm} & 
42.\ \rule{1.2cm}{0.15mm} & \\
\end{tabular}



\vspace{1em}

\textbf{MULTIPLE CHOICE.} Choose the one alternative that best completes the statement or answers the question (1 point each, total of 45 points).

\vspace{1em}

% --------------------------------------------------------------------
% Multiple Choice Questions
% --------------------------------------------------------------------

\begin{enumerate}[label=\textbf{\arabic*.}]

\item The slope of a country's production possibility frontier with cloth measured on the horizontal and food measured on the vertical axis in the Ricardian model is equal to \underline{\hspace{2em}} and it \underline{\hspace{2em}} as more cloth is produced.
\begin{enumerate}[label=\Alph*.]
\item $-MPL_F/MPL_C$; becomes flatter
\item $-MPL_F/MPL_C$; becomes steeper
\item $-MPL_F/MPL_C$; is constant
\item $-MPL_C/MPL_F$; is constant
\item $-MPL_C/MPL_F$; becomes steeper
\end{enumerate}

\item Economists consider the effects of free trade on income distribution to be \underline{\hspace{2em}} important than the effects on overall welfare because \underline{\hspace{2em}}.
\begin{enumerate}[label=\Alph*.]
\item less; those who are harmed can be compensated by those who gain
\item more; the effects on income distribution are major and consequential
\item more; those who are harmed are not compensated by those who gain
\item less; the effects on income distribution are minor and inconsequential
\item less; the wealthy benefit and only the poor lose
\end{enumerate}

\item In the specific factors model, which of the following will increase the quantity of labor used in cloth production?
\begin{enumerate}[label=\Alph*.]
\item a decrease in the price of labor
\item an increase in the price of food relative to that of cloth
\item an increase in the price of cloth relative to that of food
\item an equal percentage increase in the price of food and cloth
\item an equal percentage decrease in the price of food and cloth
\end{enumerate}

\item In the specific factors model, labor is defined as a(an)
\begin{enumerate}[label=\Alph*.]
\item intensive factor.
\item mobile factor.
\item specific factor.
\item fixed factor.
\item variable factor.
\end{enumerate}

\item The specific factors model assumes that there are \underline{\hspace{2em}} goods and \underline{\hspace{2em}} factor(s) of production.
\begin{enumerate}[label=\Alph*.]
\item four; three
\item three; two
\item two; one
\item two; two
\item two; three
\end{enumerate}

\item In the specific factors model, the effects of trade on welfare are \underline{\hspace{2em}} for mobile factors, \underline{\hspace{2em}} for fixed factors used to produce the exported good, and \underline{\hspace{2em}} for fixed factors used to produce the imported good.
\begin{enumerate}[label=\Alph*.]
\item positive; ambiguous; ambiguous
\item ambiguous; positive; negative
\item positive; positive; positive
\item ambiguous; negative; positive
\item negative; ambiguous; ambiguous
\end{enumerate}

\item Factors tend to be specific to certain uses and products
\begin{enumerate}[label=\Alph*.]
\item in the short run.
\item in capital-intensive industries.
\item in labor-intensive industries.
\item in countries lacking comparative advantage.
\item in countries lacking fair labor laws.
\end{enumerate}

\item According to the Heckscher--Ohlin model, the source of comparative advantage is a country's
\begin{enumerate}[label=\Alph*.]
\item technology.
\item factor endowments.
\item human capital.
\item advertising.
\item political system.
\end{enumerate}

\newpage
\item According to the Heckscher--Ohlin model
\begin{enumerate}[label=\Alph*.]
\item everyone gains from trade.
\item a country gains from trade if its exports have a high value added.
\item the gainers from trade could compensate the losers and still retain gains.
\item only the country with the more advanced technology gains from trade.
\item the scarce factor gains from trade and the abundant factor loses.
\end{enumerate}
\vspace{0.75 cm}
\item Biased expansion of production possibilities occurs
\begin{enumerate}[label=\Alph*.]
\item when the production possibility frontier shifts outward equally in the production of all goods.
\item when the production possibility frontier shifts inward more in one direction than in the other.
\item when the production possibility frontier shifts inward equally in the production of all goods.
\item when the production possibility frontier shifts outward much more in one direction than in the other.
\end{enumerate}
\vspace{0.75 cm}
\item Suppose that there are two factors, capital and land, and that the United States is relatively land-endowed while the European Union is relatively capital-endowed. According to the Heckscher--Ohlin model
\begin{enumerate}[label=\Alph*.]
\item all landowners should support free trade.
\item all capitalists in both countries should support free trade.
\item the U.S.\ should compensate European countries once trade commences.
\item European capitalists should support U.S.--European free trade.
\item European landowners should support U.S.--European free trade.
\end{enumerate}
\newpage

\item Which of the following is an assertion of the Heckscher--Ohlin model?
\begin{enumerate}[label=\Alph*.]
\item Factor endowments determine the technology that is available to a country, which determines the good in which the country will have a comparative advantage.
\item An increase in a country's labor supply will increase production of both the capital-intensive and the labor-intensive good.
\item In the long run, labor is mobile and capital is not.
\item Factor price equalization will occur only if there is costless mobility of all factors across borders.
\item An increase in a country's labor supply will increase production of the labor-intensive good and decrease production of the capital-intensive good.
\end{enumerate}

\item An increase in the real interest rate, all other things held constant, will cause a country's \underline{\hspace{2em}} to \underline{\hspace{2em}}.
\begin{enumerate}[label=\Alph*.]
\item terms of trade; improve
\item terms of trade; worsen
\item welfare level; improve
\item current consumption; increase
\item current consumption; decrease
\end{enumerate}

\item An increase in a country's net commodity terms of trade will
\begin{enumerate}[label=\Alph*.]
\item not always guarantee positive changes in the country's economy.
\item always increase the country's real income.
\item always increase the country's production of its import-competing good.
\item always increase the country's economic welfare.
\item never increase the country's quantity of exports.
\end{enumerate}

\item The distinctive feature of tariffs and export subsidies is that
\begin{enumerate}[label=\Alph*.]
\item they make imported and exported goods cheaper.
\item they give incentives for importers and exporters.
\item they create a difference between prices at which goods are traded on the world market and prices at which those goods can be purchased within a country.
\item they equalize prices at which goods are traded on the world market and prices at which those goods can be purchased within a country.
\end{enumerate}
\newpage
\item A rise in the relative price of cloth leads the economy to produce \underline{\hspace{2em}} cloth and \underline{\hspace{2em}} food, therefore \underline{\hspace{2em}} the relative supply of cloth.
\begin{enumerate}[label=\Alph*.]
\item more; less; decrease
\item more; less; increase
\item less; more; increase
\item less; more; decrease
\end{enumerate}

\item The effect of an export subsidy is to
\begin{enumerate}[label=\Alph*.]
\item give producers an incentive to export.
\item make it less profitable to sell abroad than at home.
\item raise government's revenue.
\item raise the prices of exported goods outside a country.
\end{enumerate}

\item If points A and B are two locations on a country's production possibility frontier, then
\begin{enumerate}[label=\Alph*.]
\item consumers are indifferent between the two bundles.
\item the country could produce either of the two bundles.
\item both bundles must have the same relative cost.
\item at any point in time, the country could produce both.
\item producers are indifferent between the two bundles.
\end{enumerate}

\item If the ratio of the price of cloth ($P_C$) divided by the price of food ($P_F$) increases in the international marketplace, then
\begin{enumerate}[label=\Alph*.]
\item the terms of trade of all countries will improve.
\item the terms of trade of cloth exporters will worsen.
\item the terms of trade of cloth exporters will improve.
\item all countries would be better off.
\item the terms of trade of food exporters will improve.
\end{enumerate}

\newpage
\item Patterns of interregional trade are primarily determined by \underline{\hspace{2em}} rather than \underline{\hspace{2em}} because factors of production are generally \underline{\hspace{2em}}.
\begin{enumerate}[label=\Alph*.]
\item external economies; population; immobile
\item population; external economies; immobile
\item internal economies; population; immobile
\item external economies; natural resources; mobile
\item internal economies; external economies; mobile
\end{enumerate}

\item If a firm's output more than doubles when all inputs are doubled, production is said to occur under conditions of
\begin{enumerate}[label=\Alph*.]
\item decreasing returns to scale.
\item constant returns to scale.
\item increasing returns to scale.
\item imperfect competition.
\item intra-industry equilibrium.
\end{enumerate}

\item If some industries exhibit internal increasing returns to scale in each country, we should not expect to see
\begin{enumerate}[label=\Alph*.]
\item intra-industry trade between countries.
\item inter-industry trade between countries.
\item increased productivity in both countries.
\item high levels of specialization in both countries.
\item perfect competition in these industries.
\end{enumerate}

\item External economies of scale often arise because similar firms
\begin{enumerate}[label=\Alph*.]
\item agree to cooperate to expand global trade.
\item have excellent internal logistics.
\item have economies of scale in production.
\item locate in the same geographic region.
\item collude to fix prices and increase profits.
\end{enumerate}
\newpage
\item External economies of scale will \underline{\hspace{2em}} cost per unit when output is \underline{\hspace{2em}} by \underline{\hspace{2em}}.
\begin{enumerate}[label=\Alph*.]
\item increase; increased; the industry
\item reduce; reduce; the industry
\item reduce; increased; a firm
\item reduce; increased; the industry
\item increase; increased; a firm
\end{enumerate}

\item If output is increased in the long run, average production costs in the presence of internal economies of scale will \underline{\hspace{2em}}, and in the presence of external economies of scale, will \underline{\hspace{2em}}.
\begin{enumerate}[label=\Alph*.]
\item increase; decrease
\item decrease; remain constant
\item decrease; decrease
\item increase; remain constant
\item remain constant; increase
\end{enumerate}

\item The imposition of tariffs on imports results in deadweight (triangle) losses. These are
\begin{enumerate}[label=\Alph*.]
\item production and consumption distortion effects.
\item distortion of incentives.
\item revenue effects.
\item efficiency effects.
\item redistribution effects.
\end{enumerate}

\item The principle benefit of tariff protection goes to
\begin{enumerate}[label=\Alph*.]
\item foreign consumers of the good produced.
\item domestic producers of the good produced.
\item foreign producers of the good produced.
\item domestic consumers of the good produced.
\item the domestic government.
\end{enumerate}
\newpage
\item An export subsidy \underline{\hspace{2em}} prices in the exporting country while \underline{\hspace{2em}} them in the importing country.
\begin{enumerate}[label=\Alph*.]
\item lowers; lowering
\item lowers; raising
\item raises; lowering
\item raises; raising
\end{enumerate}

\item Tariff rates on products imported into the U.S.
\begin{enumerate}[label=\Alph*.]
\item were the government's main source of income in 2006.
\item reached an all time high in 2002.
\item were prohibited by the Constitution.
\item have dropped substantially over the past 50 years.
\item have risen steadily since 1920.
\end{enumerate}

\item Trade between British Columbia and a Canada province is likely much larger than trade with an equally distant U.S. state because of
\begin{enumerate}[label=\Alph*.]
\item size of local economy.
\item religion.
\item population.
\item negative effects of the borders.
\end{enumerate}

\item A situation where a final product goes through many production stages by different companies in different countries instead of being produced by one company in one country is called
\begin{enumerate}[label=\Alph*.]
\item horizontal integration.
\item vertical disintegration.
\item vertical integration.
\item horizontal disintegration.
\end{enumerate}
\newpage
\item Since the early 1970s, world's trade as a share of world production has
\begin{enumerate}[label=\Alph*.]
\item increased slightly before dropping off.
\item increased.
\item decreased.
\item remained constant.
\item fluctuated widely with no clear trend.
\end{enumerate}

\item In the present, most of the exports from China are
\begin{enumerate}[label=\Alph*.]
\item technology-intensive products.
\item services.
\item primary products including agricultural.
\item manufactured goods.
\item overpriced by world market standards.
\end{enumerate}

\item In general, which of the following do \emph{not} tend to increase trade between two countries?
\begin{enumerate}[label=\Alph*.]
\item larger economies
\item linguistic and/or cultural affinity
\item mutual membership in preferential trade agreements
\item historical ties
\item the existence of well controlled borders between countries
\end{enumerate}

\item Which of the following statements is TRUE?
\begin{enumerate}[label=\Alph*.]
\item The competitive advantage of an industry depends not only on its productivity relative to the foreign industry, but also on the domestic wage rate relative to the foreign wage rate.
\item The competitive advantage of an industry only depends on the domestic wage rate relative to the foreign wage rate.
\item The competitive advantage of an industry depends on a country's absolute productivity advantage over other countries.
\item The competitive advantage of an industry only depends on its productivity relative to the foreign industry.
\end{enumerate}

\newpage
\item When one country can produce a unit of a good with less labor than another country, the first country has a(n) \underline{\hspace{2em}} in producing that good.
\begin{enumerate}[label=\Alph*.]
\item advanced advantage
\item comparative advantage
\item absolute advantage
\item comparable advantage
\end{enumerate}

\item Assume $a_{LW}$ and $a_{LC}$ are the unit labor requirements in wine and cheese production, respectively. The opportunity cost of cheese in terms of wine is
\begin{enumerate}[label=\Alph*.]
\item $a_{LW}/a_{LC}$
\item $1/a_{LC}$
\item $a_{LC}/a_{LW}$
\item $1/a_{LW}$
\end{enumerate}

\item A country engaging in trade according to the principles of comparative advantage gains from trade because it
\begin{enumerate}[label=\Alph*.]
\item is producing exports indirectly more efficiently than it could alternatively.
\item is producing imports indirectly using fewer labor units.
\item is producing imports indirectly more efficiently than it could domestically.
\item is producing exports using fewer labor units.
\item is producing exports while outsourcing services.
\end{enumerate}

\item Which of the following statements is TRUE?
\begin{enumerate}[label=\Alph*.]
\item A country's exports are determined by its relative high productivity in the industry compared with relative productivity in other sectors.
\item A country's exports are determined by its high productivity in the industry compared with that of foreigners.
\item A country's exports are determined by its ability to match other countries' productivity.
\item A country's exports are determined by its low productivity in the industry compared with that of foreigners.
\end{enumerate}

\newpage
\item The country of Rhozundia is blessed with rich copper deposits. The cost of copper produced (relative to the cost of widgets produced) is therefore very low. From this information we know that
\begin{enumerate}[label=\Alph*.]
\item Rhozundia has a comparative advantage in copper.
\item Rhozundia may or may not have a comparative advantage in copper.
\item Rhozundia should export both widgets and copper.
\item Rhozundia should invest in more widget production.
\item Rhozundia should import copper and export widgets.
\end{enumerate}

\item When a country can produce a good or service at a lower opportunity cost than another country, the country is said to have a(n)
\begin{enumerate}[label=\Alph*.]
\item absolute advantage
\item comparative advantage
\item cost advantage
\item trading advantage
\end{enumerate}

\item Which of the following statements is TRUE?
\begin{enumerate}[label=\Alph*.]
\item Migrating workers often have very similar characteristics as native workers in the receiving country.
\item Migrating workers often have very different characteristics than native workers in the receiving country.
\item Migrating workers often underperform native workers in the receiving country.
\item Migrating workers often outperform native workers in the receiving country.
\end{enumerate}

\item The slope of a country's production possibility frontier is equal to \underline{\hspace{2em}} and the optimal production point is located where the slope is equal to \underline{\hspace{2em}}. Assume that output of good $Y$ is measured on the vertical axis, output of good $X$ is measured on the horizontal axis, $MPL$ is the marginal product of labor with a subscript indicating which good, $P$ is the price of a good, and $w$ is the wage rate.
\begin{enumerate}[label=\Alph*.]
\item $-P_X/P_Y$; $-MPL_Y/MPL_X$
\item $-MPL_Y/MPL_X$; $-P_X/P_Y$
\item $-MPL_Y/w$; $-MPL_F/w$
\item $-P_X/w$; $-P_Y/w$
\item $-MPL_X/MPL_Y$; $-P_X/P_Y$
\end{enumerate}

\item Refer to the figure above, which shows a country's possible production possibility frontiers and indifference curves. If the country is producing at \underline{\hspace{2em}}, then moving to \underline{\hspace{2em}} will cause utility to \underline{\hspace{2em}}.
\begin{enumerate}[label=\Alph*.]
\item point a; point b; increase
\item point c; point b; increase
\item point a; point c; remain unchanged
\item point c; point b; decrease
\item point c; point b; remain unchanged
\end{enumerate}

% (Insert the relevant PPF/indifference-curve figure here.)

\item Given the information in the table above, Foreign's opportunity cost of cloth is
\begin{enumerate}[label=\Alph*.]
\item 0.5
\item 2.0
\item 6.0
\item 1.5
\item 3.0
\end{enumerate}

% (Insert the relevant table for Foreign's technologies/output here.)

\end{enumerate}


\newpage

\section*{Problem 1}

China and Mexico can both produce drones. Average cost in each country depends on domestic output because of external economies of scale:
\[
AC_{China} = 40 - 0.02Q_{China}, \qquad AC_{Mexico} = 42 - 0.01Q_{Mexico},
\]
where $AC_i$ is the average cost in country $i$ and $Q_i$ is domestic output.

World demand for drones is:
\[
Q^{World} = 300.
\]

Initially, China produces $Q_{China} = 250$ and Mexico produces $Q_{Mexico} = 50$.

\begin{enumerate}[label=(\arabic*)]
\item Compute the average cost in each country and identify the lower-cost producer.
\item Suppose Mexico’s domestic market expands, raising its drone production to $Q_{Mexico} = 200$. Recompute the average costs for China and Mexico.
\item Based on these revised costs, determine which country is more likely to become the world’s drone production hub.
\item Explain how this example illustrates:
\begin{itemize}
    \item path dependence,
    \item multiple possible long-run equilibria,
    \item the role of expectations in determining industry location when external economies of scale are present.
\end{itemize}
\end{enumerate}


\newpage
\section*{Problem 2 }

The United States imports bicycles and faces the domestic demand and supply:
\[
Q_D = 1{,}200 - 3P, \qquad Q_S = 2P - 100.
\]

The world price of bicycles is $P_W = 200$.  
The United States is a large importer, and foreign export supply is:
\[
P^{Foreign} = 150 + 0.25Q^{Foreign}.
\]

The U.S.\ imposes an import quota of 120 bicycles.

\begin{enumerate}[label=(\arabic*)]
\item Find the free-trade domestic price in the U.S.\ and the free-trade quantity of imports.
\item Under the quota, determine:
\begin{itemize}
    \item the domestic price $P_Q$,
    \item the foreign export price $P^{Foreign}$,
    \item the value of quota rents.
\end{itemize}
\item Suppose quota licenses are given to U.S.\ firms. Compute:
\begin{itemize}
    \item the production distortion,
    \item the consumption distortion,
    \item quota rents,
    \item the terms-of-trade effect.
\end{itemize}
\item Compare these welfare effects with a tariff that restricts imports to the same quantity. Which policy is more efficient, and why?
\end{enumerate}

\end{document}
